\documentclass[../main.tex]{subfiles}
\begin{document}
%==========================================TIÊU ĐỀ TẬP========================================
\begin{table}[h]
  \begin{adjustwidth}{-1cm}{}
    \begin{tabular}{>{\centering\arraybackslash}p{6cm}|>{\raggedright\arraybackslash}p{12.5cm}}
      \multirow{5}{6cm}
      % Cột bên trái: ÉPISODE và số tập
      {\\[-40pt]
        \flaregothic\fontsize{20pt}{20pt}\selectfont \centering
        {\contour{errorfour}{\textcolor{black}{ÉPISODE}}}\color{black}\\
        \burbank\fontsize{72pt}{72pt}\selectfont
        \includegraphics[height=3.12359cm]{../../../../Image/Book?/number/num18.png}
        \\[18pt]
        % \flaregothic\fontsize{14pt}{14pt}\selectfont
        % \color{epattr}BIENVENUE, \uline{\color{Banpire} Mel} !
        % \flaregothic\fontsize{18pt}{18pt}\selectfont
        % \color{epattr}ÉPISODE PERDU
        \color{black}
      }
       &
      % Dòng thứ nhất: tên tập tiếng Nhật
      {\fotskip\fontsize{18pt}{18pt}\selectfont \ruby{森}{もり}の\ruby{記}{き}\ruby{憶}{おく}}
      \\[6pt]
      % Dòng thứ hai: tên tập tiếng Anh
       & {\cafeteria\fontsize{18pt}{18pt}\selectfont Memories of the Woods}                \\[4pt]
      % Dòng thứ ba: tên tập tiếng Việt
       & {\vnmsans\fontsize{18pt}{18pt}\selectfont Ký ức từ khu rừng}                      \\[8pt]
      % Dòng thứ tư: Nhân vật xuất hiện
       & {\caviar{\fontsize{14pt}{14pt}\selectfont Track used: The Mystery - ChronosAeon}}
      \\[6pt]
      % Dòng cuối cùng: Thời gian diễn ra của tập (thường sẽ là ngày phát hành)
       & {\continuum\fontsize{16pt}{16pt}\selectfont Ngày phát hành: 18/06/2022}           \\[8pt]
    \end{tabular}
  \end{adjustwidth}
\end{table}
%===========================================NỘI DUNG==========================================
\color{black}

Khi ánh sáng mờ dần, họ đứng giữa một không gian trôi nổi. Mọi thứ quanh họ đều bị phủ một lớp sương mỏng màu bạc. Từng nét vẽ, từng mảng màu chập chờn giữa tồn tại và hư ảo.

``Hình ảnh này\dots\ không phải là quá khứ thông thường,'' Game bất chợt lên tiếng. ``Có lẽ đây là mảnh ký ức được bức tranh lưu giữ.''

Shino nhìn xuống bàn tay mình. Nó trong suốt, như thể chỉ còn là một nét mực run rẩy trên nền vẽ. Cô đưa tay định chạm vào một bông hoa đang lơ lửng, nhưng ngón cô đi xuyên qua như đâm vào làn khói.

``Chúng ta\dots\ đang ở trong tranh sao?''

Kuon khẽ gật đầu, ánh mắt cậu vẫn dõi theo những đường cong đang biến hóa quanh họ. ``Đúng hơn\dots\ là trong trí nhớ của người đã tạo ra nó.''

Không gian chớp sáng. Những đường cọ ảo ảnh bỗng chuyển động, xếp lại thành hình cây, đất, và bầu trời. Màu xanh lan rộng như mực loang trong nước. Tiếng gió, tiếng chim, mùi nhựa cây, tất cả cùng ùa về.

\img{e3-2a}

Một khu rừng hiện ra, như thể vừa được vẽ ra từ những nét mực còn ướt. Cây cối chen chúc, thân cây xoắn ốc vút lên cao, vỏ nứt nẻ đan xen thành những hình mắt xích nhìn chằm chằm xuống đất. Tán lá dày đặc đan thành mái vòm xanh thẫm, chỉ để lọt qua những mảng sáng mỏng như lưỡi dao khẽ lướt trên mặt đất. Rêu phủ kín gốc, mềm mại như lớp da đang thở, mỗi bước chân đạp xuống liền nhúng vào tiếng *ỤP* khe khẽ, như có ai đó thở dưới lớp lá mục. Không khí nặng mùi nhựa cây và đất ẩm, hòa quyện thành một mùi ngọt ngào đến nghẹt thở, khiến mọi âm thanh dường như bị nuốt chửng chỉ còn tiếng tim đập rơi rớt trong lồng ngực. Một tia sáng le lói rọi qua kẽ lá, biến thành bụi vàng lấp lánh lơ lửng, lặng lẽ quay tròn như những con thoi đang dệt lại ký ức cũ. Khu rừng không chỉ im lặng, mà nó đang lắng nghe, đang chờ, đang ghi khắc từng nhịp thở lạ lẫm của những vị khách vừa bước vào thế giới mực và giấc mơ này.

Trong khung cảnh ấy, một cô gái dần xuất hiện, mái tóc đỏ cột thành hai bím dập dờn trong gió, đồng phục trường Aogami thời ngày xưa, trên tay là cuốn sổ ký họa. Cô bước chậm rãi, tay lật giở từng trang giấy, miệng khe khẽ ngân nga một giai điệu lạ. Nắng xuyên qua tán lá, rải thành từng vảy sáng trên tóc cô.

\img{e3-2b}

``Khung cảnh ở đây yên bình ghê\~{}'' cô gái ấy nói, hít lấy một hơi dài.

Những bước chân nhẹ nhàng của cô để lại dấu chân mờ trên thảm lá khô, tiếng lá xào xạc như đang thì thầm kể chuyện. Một chú chim nhỏ từ cành cao nhảy xuống đậu trên vai cô, líu lo vài nốt nhạc rồi bay đi, để lại không gian lại chìm trong thinh lặng.

``Cô gái đó\dots'' Kaigou khẽ nói, ``Cô ấy\dots\ hay đi dạo trong rừng à?''

Shino nhìn theo bóng dáng ấy, đôi mắt ánh lên vẻ trầm tư: ``Hình như mình nhớ\dots\ Nanase-san cũng từng làm vậy. Cảm giác y hệt.''

Kuon nhìn gương mặt đang hớn hở kia, nghiêng đầu: ``\dots Sao anh thấy trông em ấy hao hao giống Marine-senchou vậy nhỉ\dots?''

``Một tài năng khác bên công ty cậu sao, Kuon-user?''

``Ừ.''

Một làn gió thoảng qua, kéo theo tiếng lật giấy mỏng. Mari dừng lại dưới một gốc cây, mở sổ ra, bút chì bắt đầu chuyển động. Mỗi nét vẽ của cô tựa như đang sinh ra ánh sáng.

Game lên tiếng, giọng lạnh và sắc:

``Hay là cô ấy đang tìm cảm hứng. Những người bị ám ảnh bởi sáng tạo thường trốn vào nơi cô lập nhất.''

Kaigou nhìn quanh, khẽ nhíu mày: ``Một thế giới được tái hiện từ ký ức\dots\ nghĩa là mọi thứ ở đây đều có ý nghĩa. Có lẽ ta nên xem cô ấy định vẽ gì.''

Kuon khẽ gật, giọng cậu trầm lại: ``Phải. Biết đâu ta có thể tìm được manh mối.''

Cả bốn người im lặng bước theo, dõi theo bóng Mari giữa khu rừng ánh sáng, nơi từng nét cọ bắt đầu chuyển động, báo hiệu cho thứ gì đó sắp được vẽ ra, và có lẽ, cũng sắp được thức tỉnh.

\ct

Mari ngồi xuống một khúc gỗ phủ rêu. Trên đầu, những tia nắng muộn rọi qua tán cây, trải lên trang giấy trong cuốn sổ ký họa. Cô nhìn đăm đăm vào đó một lúc lâu, ngón tay khẽ run khi đặt bút chì lên giấy.

Vài nét phác thảo đầu tiên xuất hiện, đường cong của một vật thể, những mảng sáng tối đan xen, nhưng càng vẽ, ánh mắt cô càng trở nên lạc lõng.

``Không\dots\ không được. Cứ như thế này, chẳng có gì thay đổi cả\dots''

Giọng của cô nhỏ đến mức chỉ mình cô nghe thấy. Cô thở dài, nghiêng đầu, ngắm nghía bức phác thảo, rồi lại xóa đi, vẽ lại, từng nét bút trở nên chậm chạp, nặng nề như kéo theo cả tâm trí.

Shino đứng gần đó, khẽ thì thầm: ``Cô ấy\dots\ trông thật buồn.''

Suzuki nhìn cô nàng tóc đỏ kia, thở dài một hơi như muốn sẻ chia sự cảm thông.

Kaigou gật nhẹ, ánh mắt dõi theo bàn tay đang run của Mari. ``Cảm giác như cô ấy đang bị chính bức tranh phản kháng vậy.''

Kuon nói chậm, giọng anh khàn đi giữa tiếng gió xào xạc: ``Đó là sự bế tắc mà ai từng sáng tạo đều hiểu\dots\ Khoảnh khắc mà mọi ý tưởng tưởng rằng đều đã ở trong tầm tay, bỗng chốc trở nên xa vời tới lạ.''

Mari đặt bút xuống, chống cằm lên tay, ngẩng nhìn những kẽ lá. Khuôn mặt cô đượm buồn, nhưng vẫn có chút ánh sáng trong mắt, một tia hy vọng nhỏ nhoi chưa tắt hẳn. Sau một hồi trầm ngâm, cô cất sổ vào túi, đứng dậy phủi áo.

``Thôi, hôm nay tới đây là đủ rồi.''

Cô bước đi dọc theo lối mòn dẫn ra khỏi rừng. Khi nắng dần ngả, bóng cây kéo dài trên mặt đất, Mari khẽ mỉm cười. Một nụ cười nhẹ nhõm, nhưng phảng phất đâu đó một niềm tiếc nuối nhẹ.

Ở phía xa, một bóng người xuất hiện: một cô gái với mái tóc đen buộc hai bím thấp, trên cổ tay trái là vòng tay đỏ gắn ba chiếc chuông nhỏ khẽ reo leng keng mỗi khi cô di chuyển.

Mari chớp mắt, hơi nghiêng đầu. ``Kia chẳng phải là\dots\ Reiko-san đó sao?''

Khi nhìn rõ khuôn mặt cô gái, cô nàng hai màu mắt vẫy tay, giọng vui hẳn lên: ``\textbf{Này!}''

Reiko quay lại, khẽ ngạc nhiên, rồi mỉm cười: ``À, Mari-san đó sao?''

Mari tiến lại gần, hơi cười đùa: ``Không ngờ là gặp cậu ở đây đấy. Cậu sống gần đây à?''

Reiko gật đầu, giọng nhẹ như gió: ``Ừm. Từ đây về tới nhà cũng chỉ mất vài phút thôi mà.''

Mari tròn mắt, giọng pha chút ghen tị: ``Hể\~{} Thật sao~{} Số cậu đỏ ghê. Giá mà nhà mình sống ở gần khu rừng này, ngày nào mình cũng có thể đi dạo ở đây quá.''

Hai cô gái vừa nói chuyện vừa rảo bước về phía con đường nhỏ dẫn ra khỏi rừng, tiếng chuông nhỏ trên tay Reiko khẽ ngân đều theo từng nhịp chân.

Từ phía xa, Suzuki bỗng khựng lại. Cô lùi nửa bước, mắt mở to khi nhìn thấy hình ảnh cô gái đeo vòng tay đỏ ấy. Mái tóc buộc hai bím thấp, đôi mắt hai màu - đỏ và xanh dương - ánh lên giữa hoàng hôn, và cả chiếc vòng tay màu đỏ với ba chiếc chuông, thứ mà cô vẫn đang đeo trên tay.

``Đó là\dots\ mình ư?''

Không khí chùng xuống. Kaigou quay sang, giọng cô nhỏ và chậm: ``\dots Dù chị không muốn nói, nhưng thực sự đó là em đó, Suzu. Thời điểm khi em còn chưa nhớ được bản thân mình là ai.''

Kuon im lặng một thoáng, rồi khẽ nghiêng đầu, nói: ``Nếu thế thì bức tranh này không chỉ chứa ký ức của Mari\dots\ mà còn lưu lại một phần ký ức của em nữa, Suzuki.''

Suzuki cúi đầu, bàn tay siết nhẹ cổ tay mình như cố kiểm tra xem chiếc vòng đỏ kia có thật hay không.

``Một phần\dots\ ký ức của em sao\dots?''

Mari và Reiko vẫn tiếp tục đi, giọng nói của họ vang vọng như tiếng vọng từ quá khứ xa xôi, dường như không để tâm gì tới sự xuất hiện của bốn con người.

Cô chị chợt hỏi khẽ: ``Suzu, em có nhớ gì về cô gái đó không? Về Mari?''

Suzuki nhắm mắt, cố gắng tìm kiếm trong trí nhớ mờ nhạt. Cô khẽ cất giọng, vừa đếm ngón tay: ``Tên đầy đủ của cô ấy là\dots\ Akagane Mari. Học cùng khối, lớp bên cạnh. Cô ấy thích vẽ, còn em thì làm thủ công.
Bọn em có từng trao đổi ý tưởng với nhau vài lần. Và một trong số đó\dots\ cô ấy bảo muốn vẽ thứ gì đó mà không ai có thể vẽ được.''

Kuon gật đầu, ánh nhìn anh dõi theo hình ảnh hai cô gái khuất dần sau màn sáng. ``Vẽ thứ mà không ai có thể vẽ được\dots\ Nghe giống một lời nguyền hơn là mục tiêu.''

Không gian như ngưng đọng. Bốn bóng người đứng giữa khu rừng ký ức, mỗi người chìm trong dòng suy tư riêng. Gió thôi thì thầm, lá thôi xào xạc; chỉ còn tiếng tim đập rơi rớt trong lồng ngực. Kaigou siết nhẹ vạt áo, nhớ tới những bức tranh mình từng vẽ rồi lại xé đi, những mảnh ký ức không tên cũng đang xé mình cô từng mảnh. Suzuki ôm lấy vai mình, mắt hai màu hằn lên vệt sáng mờ ảo của chiếc vòng đỏ ngày nào, ký ức và hiện tại chồng lấp, cô không biết mình đang nhớ hay đang quên.

Shino nhíu mày, bất chợt thắc mắc: ``Nhưng mà Suzuki-san, nếu hai người học cùng khối, sao cậu chưa từng nhắc đến cô ấy trước đây? Cả thời gian cậu học cùng với mọi người, cậu không một lần kể về Mari-san cả.''

``\dots Anh cũng thấy lạ đấy,'' Kuon nói. ``Giờ Shino-san nói anh mới để ý.''

Cô em út cúi đầu, lòng bàn tay siết lại. ``Em cũng không hiểu nữa. Có lẽ vì ký ức về cô ấy bị chôn quá sâu. Mỗi lần cố nhớ lại, dường như một thế lực nào đó ngăn cho em tiếp cận tới nó vậy.''

Kaigou phân tích với giọng sắc lẹm: ``Chị nghĩ là\dots\ khi chúng ta còn là Reiko, Kyuen và Saikyou, chúng ta chưa biết rõ tình hình mọi thứ quanh nơi này, và từ đó quên một vài mảnh ký ức. Mari có lẽ nằm trong số đó.''

Suzuki gật nhẹ, đôi mắt hai màu ánh lên vẻ buồn bã. ``Em chỉ nhớ mơ hồ có người từng đứng ngoài cửa lớp em, cầm cuốn sổ vẽ, mỉm cười. Nhưng khuôn mặt ấy bị sương mù che kín. Đến tận bây giờ, khi em biết rõ được chân tướng mọi thứ, em mới nhớ ra được ạ.''

``\dots Có vẻ như nơi này không đơn thuần làm xáo trộn ký ức của mọi người,'' Game chợt lên tiếng, giọng có chút nghiêm trọng. ``Tôi có cảm giác nó đang cố gắng giấu nhẹm chúng ta một điều gì đó. Một thứ đáng sợ hơn như thế.''

Bóng tối từ từ dâng lên, nuốt chửng lối đi nhỏ. Bốn người đứng thành hàng, không ai bước trước, như thể bất kỳ bước chân nào cũng có thể đánh thức thứ đang ngủ trong màn đêm. Shino đưa tay lên ngực, cảm nhận nhịp tim lạc nhịp, mỗi lần cô thở, không khí lạnh toát qua kẽ răng. Kuon ngẩng lên, trời đã tắt nắng, những vì sao không xuất hiện, chỉ còn khoảng trống đen ngòm treo lơ lửng.

Kaigou siết chặt nắm tay: ``Nếu đúng vậy, thì con đường mà họ đang đi\dots\ sẽ dẫn chúng ta đến câu trả lời.''

Cả bốn người im lặng, dõi theo hai bóng dáng nhỏ bé đang rời khỏi rừng, để lại phía sau là những tia sáng lấp loáng như đang bị nuốt dần vào màn đêm đang xuống.

Không gian quanh họ đột nhiên nhòe đi.

Những tán cây, con đường mòn và cả ánh sáng hoàng hôn nơi cánh rừng tan chảy như một bức tranh đang bị xóa, rồi hợp lại thành những vệt sáng đỏ-cam. Từng mảng hình ảnh xoay tròn như tua nhanh trên một cuộn phim, trước khi dừng lại trong một cú chớp sáng khô khốc.

*CẠCH*

Bốn người giật mình. Họ không còn đứng giữa rừng nữa.
Trước mắt họ, khung cảnh mới hiện ra, phòng Mỹ thuật của trường trung học Aogami, khi chiều muộn.

``Đây chẳng phải là phòng Mỹ thuật mà chúng ta vừa đứng sao?'' Suzuki thoáng ngạc nhiên.

``Trông nó còn gọn gàng hơn lúc chúng ta đứng nữa\dots'' Shino lẩm bẩm, mắt cô mở to.

Ánh hoàng hôn hắt qua khung cửa, nhuộm cả căn phòng bằng màu đỏ hổ phách. Không gian yên ắng đến mức có thể nghe thấy tiếng quạt trần quay chậm rãi.

Ở giữa căn phòng ấy, Mari ngồi lặng trước giá vẽ, mái tóc sẫm màu thả xuống vai, tay cô vẫn cầm chiếc cọ đã nhuốm màu sơn khô. Dưới ánh nắng cuối ngày, những lọ màu trên bàn phản chiếu ánh đỏ rực như máu.

\img{e3-2c}

``Lại một buổi chiều như thế này nữa\dots'' cô thở dài, giọng mệt mỏi.

Mari cúi đầu, khẽ nhắm mắt, trong đầu cô chợt vang lên lời nói của một người khác: một cô gái với mái tóc hồng với chiếc chuông cài trên tóc, đôi mắt xanh lá mạ như sương mù: Shinomiya Sakura.

``Mari-san, cậu có biết vì sao người ta cấm đi vào rừng phía tây không?''

``Vì\dots\ có ma à?'' cô đã trả lời đùa như thế.

``Không phải ma. Là những linh hồn bị chôn sống.''

Hình ảnh ký ức thoáng hiện lên như một giấc mơ bị vỡ vụn, lời nói của Sakura vang vọng trong căn phòng Mỹ thuật hiện tại.

``Thời xa xưa, thị trấn Aogami đã từng xảy ra nhiều trận sạt lở đất thường xuyên, nên để làm vừa lòng những vị thần núi ở nơi đây, họ đã thực hiện nghi thức hiến tế, chôn sống những đứa trẻ còn sống.''

Kuon chau mày, giọng khàn đặc: ``Chôn sống\dots\ cả trẻ con ư?''

Suzuki lùi một bước, tay siết chặt vòng chuông đỏ, mắt hai màu trống rỗng: ``N-Nghe có vẻ dã man thật\dots''

``Kể từ đó, đất núi không còn lở nữa, nhưng\dots\ những linh hồn đó vẫn còn quanh quẩn ở đâu đó, ôm mối hận của mình.''

Kaigou cắn môi, đôi môi run lẩy bẩy: ``Mê tín đến mức này thì\dots\ đáng sợ thật.''

Shino ôm vai mình, giọng thì thào: ``Chúng ta đang đứng trên nấm mồ ấy\dots\ phải không?''

``Người ta bảo, tuyệt đối không được bước vào khu rừng đó, vì khi nghe tiếng gọi của rừng, họ sẽ không bao giờ trở lại như trước. Nhưng rồi, những người dân khờ khạo vẫn bỏ mặc những lời cảnh báo đó mà tiếp tục khai thác gỗ trong những cánh rừng kia, và tiếp tục sử dụng chúng để xây dựng thị trấn Aogami đến tận ngày nay.''

Mari mở mắt. Cô khẽ rùng mình, rồi bật cười khô khốc. ``Linh hồn oán hận ư? Nghe cứ như truyện cổ tích ấy\dots''

Dù bản thân mình cố tỏ ra không tin với lời đồn đó,  trong mắt cô vẫn ánh lên một tia sợ hãi mờ nhạt. Một nỗi sợ rất mơ hồ đang chờ chực trào dâng, như thể có bàn tay vô hình đang siết lấy trái tim cô, thì thầm rằng những lời nguyền ấy chưa bao giờ chỉ là truyện cổ tích.

Ở bên ngoài, Kuon, Kaigou, Suzuki, Shino và Game lặng lẽ quan sát. Sau khi nghe những gì cô ấy nói, cả bốn người họ đều không khỏi run rẩy.

Cô chị lớn cau mày: ``Câu chuyện này\dots\ tôi nhớ rồi. Hình như Sakura-san cũng có nhắc qua về cái lời nguyền đó.''

Cô em út khẽ gật: ``Phải\dots\ hồi đó em còn nghĩ chỉ là lời đồn. Nhưng xem ra\dots\ không đơn giản thế đâu.''

Cậu con trai đan tay trước ngực, ánh mắt trầm hẳn xuống: ``Nếu điều đó là thật, thì khu rừng kia chính là nơi người ta đã hiến tế\dots\ Cũng có thể là nơi Mari đã `đón nhận' điều gì đó. Không chừng là cả Nanase-san nữa.''

Cô nàng tóc nâu ôm hai tay, giọng run nhẹ: ``Nghĩa là\dots\ bức tranh này, hay căn phòng đó, đều bị ám bởi chính những linh hồn ấy sao\dots?''

Từ chiếc đồng hồ trên cổ tay Kuon, giọng Game vang lên: chậm, lạnh và cứng rắn như kim loại: ``Không chỉ là bị ám. Có khả năng cao, cánh rừng đó chính là điểm khởi nguyên của mọi thứ. Các đoạn ký ức này không đơn thuần là hồi tưởng, mà là một phần vật chất hóa của năng lượng\dots\ thứ đã tạo nên lời nguyền này.''

``Ý ông là sao chứ?''

``Thì\dots''

Chưa để ông ta giải thích\dots

*CẠCH*

Tiếng kim loại rơi vang lên. Mari giật nảy người. Chiếc cọ vẽ tuột khỏi tay, lăn trên sàn.

``A\dots!''

Cô ôm đầu, gục xuống bàn. Hơi thở gấp gáp.

``Đầu mình\dots\ sao tự dưng lại nhức như vậy\dots''

\img{e3-2d}

Cô nắm chặt tóc, mắt mở to, đồng tử giãn ra, hơi thở dồn dập như thể trong đầu cô có hàng ngàn tiếng thì thầm đang vang lên cùng một lúc.

Bốn người đứng ngoài gần như đồng loạt khựng lại.

``Cơn đau đầu kia\dots\ giống hệt như của Shino và tôi khi bước vào căn phòng đó,'' Kaigou xanh mặt.

Kuon gật nhẹ, mắt anh ánh lên sự cảnh giác: ``Không đâu. Lần này khác. Cô ấy\dots\ trông như thể đang nghe thấy thứ gì đó trong đầu mình.''

Mari run rẩy, rồi bỗng dừng lại. Cô hít sâu, nắm lấy cọ vẽ bằng bàn tay dính đầy mồ hôi lạnh. Trong ánh chiều tà, cô mỉm cười, một nụ cười nửa mê man, nửa rực sáng.

``N-Nhưng\dots\ mình cảm nhận được rồi. Bức tranh hoàn hảo ấy\dots\ đang đến rất gần rồi.''

Giọng nói cô run run, nhưng đôi mắt đã đổi khác. Nó ánh lên một thứ rực rỡ méo mó, như người vừa chạm tay vào sự thật mà không còn đường quay lại.

\img{e3-2e1}

Kuon, Kaigou và Suzuki nhìn nhau. Không ai nói gì, nhưng ai cũng hiểu, đó là khoảnh khắc bi kịch của Mari bắt đầu.

\ct

Hoàng hôn chậm rãi buông xuống sau khung cửa phòng Mỹ thuật, ánh sáng đỏ nhạt hắt qua lớp kính mờ, kéo dài bóng dáng của Mari trên sàn gạch lấm tấm màu sơn. Không gian tĩnh lặng đến mức tiếng cọ chạm lên toan nghe rõ như kim khâu trên da thịt. Mọi thứ trông như một cảnh phim được tua chậm, rồi nhanh dần, cuốn theo dòng thời gian.

Bóng tối từ từ tràn vào, nuốt gọn những vệt sáng cuối cùng. Một lúc sau, đèn huỳnh quang trên trần bật sáng, thả xuống thứ ánh sáng vàng lạnh lẽo, khiến màu sơn dầu trên bàn như biến thành một mảng băng. Các cọ vẽ nằm ngổn ngang, in bóng thẳng đứng lên tường, những chiếc bút chì gãy vương vãi dưới gầm ghế. Gương mặt Mari bị chia đôi bởi bóng đèn, một nửa sáng, một nửa chìm trong bóng tối; chỉ có đôi mắt cô vẫn lấp lánh, phản chiếu ánh đèn như hai mảnh gương vỡ.

Và rồi, vài tiếng đã trôi qua.

Không khí trong hành lang dày đặc mùi bụi bẩn và màu phai úa của thời gian. Từng tia nắng cuối ngày lọt qua ô cửa sổ cao, nghiêng nghiêng, đổ xuống nền gạch hoa vằn vện những vệt sáng cam nhạt như lớp vải mỏng phủ lên xác một ký ức đang thoi thóp.

Tiếng bước chân của Reiko vang lên rời rạc, nhỏ bé, mỗi bước như đạp lên một lớp tro tàn của thời gian đã chết. Cô dừng lại trước cánh cửa phòng Mỹ thuật, tay đặt lên tay nắm kim loại lạnh ngắt, run nhè nhẹ. Mồ hôi trong lòng bàn tay cô bốc lên một mùi tanh nhẹ của sợ hãi. Cô thở ra một hơi dài, hơi thở run rẩy làm mờ đi một vùng không khí nhỏ trước mặt, rồi đẩy cửa.

``\dots Mari-san? Cậu có ở đó không?'' giọng cô hơi run. Sự lo lắng dâng lên khi không ai trả lời.

Cánh cửa mở ra, và trước mắt Reiko - cùng bốn người họ đang quan sát qua ký ức - là một khung cảnh khiến cả bọn chết lặng.

Giữa căn phòng, Mari ngồi thụp trước giá vẽ, thân hình nhỏ bé co rúm lại như một con nhím đã trút hết gai. Tóc cô xõa xuống che kín nửa khuôn mặt, những sợi tóc đỏ sẫm ướt đẫm dính vào nhau thành từng lọn, có chỗ đã khô cứng lại thành những mảng sơn đỏ sẫm. Dưới chân cô, sàn gỗ nâu tươm tất giờ thay chỗ cho một màu đỏ đen nhầy nhụa, loang rộng thành một vũng tròn, mép vũng loang ra thành những ngón tay mỏng, vươn tới tận chân tường như muốn bám víu vào một thứ gì đó đã rời khỏi thế giới này. Những tuýp màu nằm vương vãi, bị bóp nát, vỏ nhựa vỡ vụn cắm ngược lên trông như những chiếc răng gãy, những ruột màu đỏ tươi trào ra, trộn lẫn với máu thật thành một thứ bùn nhão nhớt, bóng loáng dưới ánh đèn.

\img{e3-2e2}

``\dots Máu kìa,'' Kaigou lẩm bẩm, đôi mắt cô trố ra.

Suzuki bước lùi nửa bước, mặt cắt không còn giọt máu. ``Không thể nào\dots\ cô ấy\dots\ đừng nói là?''

Kuon nhìn chằm chằm vào bức họa, rồi gật đầu vẻ trầm buồn. ``Đúng như em nghĩ đấy. Cô ấy đã dùng máu của mình để vẽ nên bức họa này. Dành cả mạng sống của mình.''

Chiếc đồng hồ trên tay Kuon khẽ rung. Giọng Game vang lên, trầm thấp như vọng từ nơi xa xăm: ``Vậy là ta đã thấy căn nguyên rồi\dots\ Cánh rừng, lời nguyền, và những kẻ bị nó chọn. Cô ta\dots\ có lẽ đã nghe được tiếng gọi đó.''

Không khí như bị nén lại, nặng đến mức mỗi hơi thở phải rặn ra từng chút. Reiko cảm thấy tim mình đang đập loạn xạ trong lồng ngực, như một con chim nhỏ bị nhốt trong lồng sắt, mỗi lần vỗ cánh lại đập vào thanh chắc. Cô đưa tay lên miệng, ngón tay run rẩy bịt kín môi, như thể chỉ cần một tiếng động nhỏ cũng có thể làm vỡ vụn cả không gian đang chênh vênh này.

``Mari-san\dots'' Reiko gọi, giọng cô run run, cảm tưởng như chân tay cô rụng rời trước cảnh tượng trước mắt.

Cô gái ngồi trước giá vẽ khẽ động đậy. Cánh tay cô run lên từng hồi, từng cơn run nhỏ như sóng gợn trên mặt nước tĩnh. Hơi thở cô phát ra từng chặp, nặng nề, như thể mỗi lần thở phải kéo cả một quả núi đỏ sẫm trong lồng ngực.

``Ồ\dots\ Là cậu sao, Reiko-san\dots'' giọng cô yếu ớt, pha chút run rẩy. ``Cậu đến\dots\ đúng lúc quá\dots'' Mỗi từ như một mảnh thủy tinh vỡ, rơi xuống nền, vỡ ra thành những mảnh nhỏ li ti, không thể gắn lại.

Trên bảng pha màu mà Mari đang cầm, màu đỏ không còn là màu được pha đê để vẽ nữa, mà đã trở thành một chất nhầy sền sệt, sáng bóng dưới ánh đèn, phản chiếu lại khuôn mặt cô như một tấm gương vỡ. Không còn ranh giới giữa sơn và máu, chỉ còn một thể thống nhất đỏ sẫm, dính vào từng kẽ móng tay, từng nếp gấp da. Con dao trét nằm trong tay cô, lưỡi dao còn dính một lớp sơn đỏ sệt, nhưng những giọt chất lỏng đỏ tươi vẫn tí tách rơi xuống, mỗi giọt tạo thành một vòng tròn nhỏ trên nền, lan rộng từng chút một, như thể thời gian đang chảy ra từ chính con dao.

Cô nàng Họa sĩ ngẩng đầu, chậm đến mức có thể nghe thấy tiếng cổ xương sống kêu. Khuôn mặt cô bê bết đỏ, những vệt đỏ không còn là vệt nữa mà đã trở thành một lớp mặt nạ bóng loáng, bóng đến mức có thể phản chiếu lại ánh mắt của người nhìn. Đôi mắt hai màu của cô trống rỗng như hai hố sâu không đáy, nhưng trong đáy hố ấy lại ánh lên một thứ sáng rực, một ngọn lửa đỏ cuối cùng đang cháy trong cơn mê sảng.

``Cuối cùng thì\dots'' cô thì thào, ``\dots mình đã vẽ được\dots\ một bức\dots\ mà không ai\dots\ vẽ được cả\dots''

Câu nói dở dang, không thể hoàn thiện trọn vẹn.

Đầu cô nghiêng sang một bên, cơ thể đổ gục về phía trước, tiếng *BỤP* khô khốc vang lên khi trán cô chạm nền. Bảng pha màu trượt khỏi tay, lăn đi, để lại một vệt đỏ tươi kéo dài như một dấu vết cuối cùng của sinh mạng vừa dứt.

\img{e3-2f}

``\textbf{Mari-san!}'' Reiko hét lên, quỳ sụp xuống, lay mạnh vai bạn mình. Cô quỳ sụp xuống, đầu gối đập xuống nền với một tiếng *BỊCH* đau điếng, nhưng cô không cảm thấy đau. Hai tay cô lay mạnh vai bạn mình, nhưng chỉ cảm thấy một khối thịt lạnh toát, cứng đờ, như thể đang lay một khúc gỗ đã ngâm trong máu. Không có hồi đáp. Chỉ còn tiếng thở hắt yếu ớt\dots\ rồi im bặt, im đến mức có thể nghe thấy tiếng máu đang ngừng chảy trong lồng ngực mình.

Kuon lặng người. Kaigou quay đi, siết chặt nắm tay. Suzuki nhìn trân trân, như không thể tin rằng họ đang chứng kiến một cái chết ngay trước mắt. Cả ba người họ đều hoàn toàn bất lực, không thể làm gì được với một sinh mạng đã rời đi ngay trước mắt họ.

Shino thì có lẽ là sốc hơn cả: mặt cô trắng bệch, nước mắt vô thức rơi ra, cảm tưởng như cô chưa bao giờ chứng kiến cảnh một ai đó chết thảm như vậy.

Giọng Game vang lên lần nữa, trầm và nặng nề: ``Giờ thì chúng ta hiểu vì sao bức tranh ấy tồn tại rồi. Nó không chỉ được vẽ bằng máu\dots\ mà bằng cả tính mạng của người tạo ra nó.'' Lời nói không tan trong không khí mà bám dính vào da thịt, lạnh buốt và ẩm ướt, như một lời nguyền vừa được thốt ra, vừa được khắc vào xương.

Ánh sáng trong phòng nhòe đi. Cảnh vật xung quanh bắt đầu tan biến, bị cuốn vào một cơn lốc đỏ, xoáy tròn, cuốn tất cả vào trong đó, như thể màn ký ức đang khép lại đầy phũ phàng, đóng sập lại với một tiếng không vang, chỉ là một tiếng chết chóc đơn độc trong không gian bị bóp nghẹt.

\dots

\dots

\ct

Không gian xung quanh họ như bị bóp méo, ánh sáng chập chờn rồi vụt tắt. Khi họ mở mắt ra, cả bốn đã trở lại căn phòng Mỹ thuật cũ, nơi bức tranh của Mari vẫn lặng lẽ đứng giữa sàn, phủ một lớp bụi mờ. Không còn máu, không còn hoàng hôn, chỉ còn lại mùi ẩm mốc, cảm giác lạnh buốt len lỏi vào từng khớp xương và bức họa nguyền rủa ngay trước mắt.

Suzuki là người đầu tiên lên tiếng, giọng run nhẹ: ``\dots Em vẫn chưa tin nổi là chúng ta vừa chứng kiến cái đó thật. Cô ấy\dots\ thực sự đã chết ở ngay đây, phải không\dots?''

Kaigou khoanh tay, đôi mắt ánh lên vẻ nặng nề: ``Không chỉ chết. Cô ta đã bị thứ gì đó nuốt chửng. Cách Mari hành động\dots\ không giống người đang tự sát. Giống như có thứ gì trong đầu cô ta ra lệnh phải vẽ vậy.''

Kuon im lặng khá lâu. Ánh nhìn cậu dừng lại ở bức tranh, ánh sáng mờ từ ô cửa hắt lên khuôn mặt khiến đôi mắt anh trông càng lạnh.

``\dots Cái chết của cô ấy không chỉ là bi kịch. Nó là kết quả của một sự dẫn dắt. Nếu ta để ý kỹ,'' cậu nói chậm rãi, ``Mari bắt đầu đau đầu sau khi đã đi dạo trong khu rừng ấy. Cộng hưởng với lời nguyền từ Sakura-san, chúng đã tạo cho cô ấy một cảm giác giống như\dots\ một thứ trong rừng đó đã nghe thấy, rồi đánh thức điều gì đó trong cô.''

Shino thoáng tái mặt đi. ``Chính xác là như thế\dots\ Cách Mari-san ôm đầu, rồi nói năng loạn trí\dots\ nó giống hệt Nanase-san.''

Cô nàng tóc nâu bỗng chốc cứng đờ, mắt mở to nhìn vào khoảng không trước mặt. Một cơn ớn lạnh chạy dọc sống lưng: những hình ảnh đêm đó lại tràn về: Nanase gào thét, cảnh các bạn học la hét khi ngôi trường đổ sập, và cả việc đích thân cô nàng tóc xanh kia công khai thừa nhận đã bắt nạt cô. Shino run nhẹ, thì thào: ``Sao mình vẫn nhớ rõ mọi thứ\dots\ dù mình không hề ở đó?''

Game vang lên từ đồng hồ, giọng có gì đó thích thú: ``Cô vẫn chưa hiểu sao, Shino-san? Tôi đã nói rồi, cô và Sora từng là một. Mọi ký ức cô ấy trải qua, dù bị chia cắt, vẫn đọng lại trong cô như bóng phản chiếu. Cô nhớ vì cô đã từng là người trong cuộc.''

Shino nuốt khan, lòng bàn tay ẩm ướt. ``Vậy là\dots\ những gì đã trải qua trong quá khứ lúc đó\dots\ tôi cũng cảm nhận được như chính mình bị tấn công?''

``Đúng thế,'' Game đáp. ``Ký ức không biết ranh giới 'tôi' và 'cô ấy'. Khi cô nhìn thấy Nanase hay Mari, cô đang thấy lại chính nỗi đau từng đè nén.''

Shino chớp mắt, tim cô như bị bóp nghẹt. Một luồng hơi lạnh chạy dọc sống lưng, không phải vì gió, mà vì những hình ảnh vụt qua: bàn tay run rẩy cầm cọ, đôi mắt đỏ ửng khi bức tranh bị xé, tiếng thở dài chìm trong họa phòng lúc chiều tàn. Cô không hề chứng kiến những khoảnh khắc ấy, thế mà mỗi khung hình lại quá đỗi thân thuộc, như thể chính cô từng đứng đó, từng thấm mùi dầu hoà loang lổ trên nền vải.

``Ký ức của cô ấy\dots\ mình cũng cảm nhận được,'' cô thì thầm, giọng khàn vì xúc động. ``Mình có thể cảm thấy nỗi cô đơn của Mari, giống hệt như lúc mình đứng cầu nguyện, không biết nên làm gì để chứng minh mình tồn tại.'' Cô ôm lấy ngực, như sợ trái tim đang đập loạn nhịp kia sẽ rơi ra ngoài. ``Chúng ta chia sẻ cùng một nỗi sợ\dots\ sợ bị lãng quên.''

Kaigou nhìn sang: ``Nanase-san\dots\ là cô gái cùng lớp với cậu trước kia đúng không?''

``Phải,'' cô nàng tóc nâu khẽ nắm chặt tay áo mình, nét mặt hiện rõ một sự ám ảnh. ``Đêm thử thách gan dạ đó\dots\ cô ấy cũng như thế. Đột nhiên nổi nóng, gào lên bắt mình phải chết. Rồi sau đó\dots\ chẳng nhớ được gì. Cứ như thể một người khác đã điều khiển cậu ấy vậy.''

Kuon khẽ gật, giọng trầm xuống: ``Tóm lại thì\dots\ cả Mari và Nanase đều bị cùng một thứ tác động.''

Suzuki nuốt khan, cố nén cảm giác bất an: ``Thứ\dots\ trong cánh rừng ấy?''

``Rất có thể,'' cậu đáp. ``Căn rừng gần trường không chỉ là nơi nguyền rủa, nó có lẽ là trung tâm của tất cả. Từ Mari, đến Nanase, rồi cả những hiện tượng chúng ta thấy.''

Game lúc đó cất giọng từ chiếc đồng hồ, vang khẽ nhưng rõ ràng trong căn phòng im ắng: ``Đó có thể là một chứng cứ có thể tựa vào. Lời đồn về những đứa trẻ bị chôn sống, và những người bị phát điên\dots\ có thể không chỉ là truyền thuyết. Cánh rừng đó dường như lưu giữ ký ức và cảm xúc của những linh hồn đã chết, rồi gieo nó vào những kẻ vô tình chạm tới.''

Một cơn gió lạnh khẽ luồn qua khe cửa. Shino rùng mình, khẽ cúi xuống\dots\ rồi bất ngờ, mũi chân cô chạm phải thứ gì cứng.

``\dots Hửm?''

\img{e3-2g}

Cô cúi xuống, nhặt lên một cuốn sách cũ, bìa sờn mép, giấy đã ngả vàng. Dưới ánh đèn mờ, dòng chữ trên tem mượn hiện ra:

\begin{center}
  ``Người mượn: Shinomiya Sakura.''
\end{center}

``\dots Sakura?'' Suzuki khẽ thốt lên. ``Cô bạn với mái tóc hồng và làm việc ở đền sao?''

Kuon bước lại gần, nheo mắt nhìn phần tiêu đề mờ nhòe ở gáy sách. ``Cô ấy từng nói là có xuống thư viện để tra cứu mấy chuyện kỳ lạ quanh Aogami. Có lẽ đây là một trong số đó.''

``Thư viện trường\dots'' Shino lẩm bẩm, cảm giác rờn rợn len lên sống lưng.

``Có lẽ mọi người sẽ tìm được thứ gì đó có ích ở đó chăng?'' giọng Pháp khàn tiếp tục đưa ra lời gợi ý, nhưng vẫn không quên mỉa mai: ``Ngoài một đống giấy lộn hoặc là đống tờ báo rách nát hết hạn ra.''

Kuon liếc nhìn đồng hồ, khẽ gật đầu. ``Phải\dots\ có thể đó là nơi chứa câu trả lời tiếp theo.''

Cô nàng tóc nâu siết chặt cuốn sách trong tay, đôi mắt ánh lên vẻ kiên định xen lẫn sợ hãi. ``Nếu Sakura-san để lại manh mối ở đó\dots\ thì chúng ta phải tìm ra. Biết đâu đấy\dots''

Cả nhóm đứng giữa căn phòng im lặng. Ngoài kia, màn đêm dày đặc và những tán cây đung đưa bên ngoài cửa sổ, như đang thì thầm gọi họ tiến sâu hơn vào bí mật của nơi này.

Từ phía bên ngoài, bầu trời thị trấn Aogami đặc quánh, như thể sương đêm cũng ngập tràn nỗi u buồn từ những gì họ vừa chứng kiến. Ánh sáng mập mờ len lỏi vào từ ô cửa sổ của ngôi trường, lấp loáng theo từng bước chân của bốn người.

Không ai nói gì trong suốt quãng đường tới đó, chỉ còn tiếng gió lùa qua hàng cây và tiếng lốp giày chạm lên mặt đất ẩm lạnh.

Shino ôm chặt cuốn sách vào ngực, đôi mắt lạc đi giữa màn đêm mịt.
``\dots Thật kỳ lạ,'' cô khẽ nói, ``rõ ràng mọi thứ chỉ là ký ức, nhưng\dots\ cảm giác đó quá thật. Mình vẫn nghe thấy tiếng thở gấp của Mari-san\dots\ và mùi máu vẫn còn đọng lại.''

Kaigou đi bên cạnh, khoanh tay sau đầu, gượng cười mà không giấu nổi vẻ trầm tư. ``Thường thì người ta nói ký ức chỉ là hình ảnh, nhưng cái thứ đó\dots\ thực tế đến mức đáng sợ. Như thể chính bức tranh đang buộc chúng ta phải chứng kiến.''

``Hay đúng hơn là muốn chúng ta nhớ lại,'' Kuon đáp, giọng cậu trầm xuống. ``Bởi có thể, những gì Mari nhìn thấy\dots\ cũng chính là thứ đang theo dõi chúng ta bây giờ.''

Suzuki khẽ ngẩng đầu, nhìn lên tán cây rung nhẹ theo gió. ``Ý anh là\dots\ rừng Aogami ấy ư? Nó có thể nghe chúng ta sao?''

Cậu không trả lời, chỉ rút tay khỏi túi áo khoác, nhìn thoáng qua mặt đồng hồ nơi Game vẫn im lặng. Kim giây trên màn hình sáng yếu ớt quay chậm, ánh lên sắc đỏ mờ như máu.

``Nghe hay không nghe không quan trọng. Điều quan trọng là các cậu đã chạm vào nó rồi. Một khi đã thấy ký ức, không còn lối quay đầu.''

Câu nói đó khiến không khí càng thêm nặng nề.

Từ xa, ánh đèn của tòa nhà thư viện dần hiện ra trong màn sương. Cửa sắt lớn vẫn khép kín, lối vào chìm trong bóng tối. Mọi thứ tĩnh lặng đến bất thường, không có cả tiếng côn trùng hay lá xào xạc.

Shino dừng lại một chút, hơi thở cô tan thành khói trong không khí lạnh.
``Trông nó khác hẳn ban ngày. Mình cảm thấy\dots\ có gì đó đáng sợ đang chờ đợi ở bên trong.''

Suzuki khẽ kéo tay cô, cố mỉm cười: ``Chúng ta cũng đã đi đến đây rồi.''

Kaigou nheo mắt nhìn quanh, thận trọng: ``Không. Cô ấy nói đúng. Có thứ gì đó\dots\ đang nhìn chúng ta.''

Kuon lập tức quay lại ngay khi cô nói xong. Đằng sau họ, con đường dài hun hút, ánh đèn huỳnh quang chớp nhẹ, rồi tắt vụt trong vài giây. Một bóng hình mờ ảo hiện ra trong khoảng tối giữa hai ngọn đèn.

Đó là một cô gái trẻ, tóc dài ngang vai, mặc đồng phục học sinh của thời hiện đại, như chính ba người. Chỉ khác ở chỗ, đôi mắt cô ánh nâu trầm như đá, không phản chiếu chút ánh sáng nào.

Cô đứng đó, yên lặng, không một cử động. Mái tóc lay động nhẹ theo gió, cùng chiếc bóng dài trải trên nền đất.

``\dots Ai đó à, Onii-sama?'' cô em út lùi lại một bước, giọng nhỏ đi.

Không ai trả lời. Khi cậu cont rai nhìn lại lần nữa, khoảng tối kia đã trống rỗng. Chỉ còn lại tiếng gió rít qua hàng cây.

Kaigou khẽ chau mày: ``Tôi thề là vừa thấy có ai ở đó.''

Kuon vẫn nhìn về phía xa, ánh mắt anh lạnh và nghi hoặc. ``\dots Cô nói đúng. Quả thật là có ai đó đang theo dõi chúng ta.''

Một cơn gió mạnh bất chợt thổi qua, làm bay tung vài trang giấy từ cuốn sổ Shino đang cầm. Cô hốt hoảng cúi xuống nhặt lại, nhưng khi ngẩng đầu lên, ánh đèn cuối con đường vụt tắt hoàn toàn, chỉ còn lại ánh sáng nhạt từ tòa thư viện phía trước.

Trong bóng đêm, giữa ranh giới của thực và ảo, có cảm giác như có ai đó vẫn đang dõi theo họm không phải bằng ánh mắt của con người, mà bằng chính ký ức bị bỏ lại.

Bước chân họ tiếp tục vang lên giữa con đường trống trải, chậm rãi tiến về phía thư viện, nơi mà câu trả lời, hay có lẽ là cơn ác mộng kế tiếp, đang chờ sẵn.

\end{document}